\documentclass[12pt,a4paper]{article}

% ── Packages ──────────────────────────────────────────────────────────────────
\usepackage[a4paper, margin=2.5cm]{geometry}
\usepackage[T1]{fontenc}
\usepackage[utf8]{inputenc}
\usepackage{lmodern}
\usepackage{microtype}
\usepackage{graphicx}
\usepackage{booktabs}
\usepackage{longtable}
\usepackage{array}
\usepackage{tabularx}
\usepackage{multirow}
\usepackage{xcolor}
\usepackage{listings}
\usepackage{amsmath}
\usepackage{amssymb}
\usepackage{hyperref}
\usepackage{enumitem}
\usepackage{titlesec}
\usepackage{fancyhdr}
\usepackage{caption}
\usepackage{subcaption}
\usepackage{setspace}
\usepackage{parskip}
\usepackage{mdframed}
\usepackage{tocloft}
\usepackage{float}

% ── Colours ───────────────────────────────────────────────────────────────────
\definecolor{exactgreen}{RGB}{39,174,96}
\definecolor{approxyellow}{RGB}{211,84,0}
\definecolor{notred}{RGB}{192,57,43}
\definecolor{partialblue}{RGB}{41,128,185}
\definecolor{codebg}{RGB}{245,245,245}
\definecolor{codeframe}{RGB}{200,200,200}
\definecolor{linkblue}{RGB}{26,115,232}
\definecolor{sectionblue}{RGB}{23,55,94}
\definecolor{lightgray}{RGB}{230,230,230}

% ── Hyperref ──────────────────────────────────────────────────────────────────
\hypersetup{
  colorlinks=true,
  linkcolor=linkblue,
  urlcolor=linkblue,
  citecolor=linkblue,
  pdftitle={Project Report: Online Grooming Detection System},
  pdfauthor={BTP-demo},
  bookmarks=true,
  bookmarksnumbered=true
}

% ── Section styling ───────────────────────────────────────────────────────────
\titleformat{\section}{\large\bfseries\color{sectionblue}}{\thesection}{1em}{}[\titlerule]
\titleformat{\subsection}{\normalsize\bfseries\color{sectionblue}}{\thesubsection}{1em}{}
\titleformat{\subsubsection}{\normalsize\itshape\color{sectionblue}}{\thesubsubsection}{1em}{}

% ── Header / footer ───────────────────────────────────────────────────────────
\pagestyle{fancy}
\fancyhf{}
\fancyhead[L]{\small\textit{Online Grooming Detection System}}
\fancyhead[R]{\small\textit{BTP-demo}}
\fancyfoot[C]{\small\thepage}
\renewcommand{\headrulewidth}{0.4pt}

% ── Code listings ─────────────────────────────────────────────────────────────
\lstdefinestyle{codestyle}{
  backgroundcolor=\color{codebg},
  frame=single,
  framerule=0.5pt,
  rulecolor=\color{codeframe},
  basicstyle=\ttfamily\small,
  breaklines=true,
  breakatwhitespace=true,
  columns=fullflexible,
  keepspaces=true,
  showstringspaces=false,
  numbers=none,
}
\lstset{style=codestyle}

% ── Coloured status badges ────────────────────────────────────────────────────
\newcommand{\exactbadge}{\textcolor{exactgreen}{\textbf{[EXACT]}}}
\newcommand{\approxbadge}{\textcolor{approxyellow}{\textbf{[APPROX.]}}}
\newcommand{\partialbadge}{\textcolor{partialblue}{\textbf{[PARTIAL]}}}
\newcommand{\notbadge}{\textcolor{notred}{\textbf{[NOT POSSIBLE]}}}

% ── Caution box ───────────────────────────────────────────────────────────────
\newmdenv[
  linecolor=approxyellow,
  backgroundcolor=yellow!8,
  linewidth=1.5pt,
  innerleftmargin=10pt,
  innerrightmargin=10pt,
  innertopmargin=8pt,
  innerbottommargin=8pt
]{cautionbox}

\newmdenv[
  linecolor=exactgreen,
  backgroundcolor=green!6,
  linewidth=1.5pt,
  innerleftmargin=10pt,
  innerrightmargin=10pt,
  innertopmargin=8pt,
  innerbottommargin=8pt
]{infobox}

% ── TOC style ─────────────────────────────────────────────────────────────────
\renewcommand{\cftsecleader}{\cftdotfill{\cftdotsep}}

% ─────────────────────────────────────────────────────────────────────────────
\begin{document}

% ══ Title page ════════════════════════════════════════════════════════════════
\begin{titlepage}
  \centering
  \vspace*{2cm}

  {\Huge\bfseries\color{sectionblue} Online Grooming Detection System \par}
  \vspace{0.4cm}
  {\Large\bfseries Project Report \par}
  \vspace{1.2cm}

  \rule{\linewidth}{0.5pt}
  \vspace{0.4cm}

  {\large\itshape Implementation of:\par}
  \vspace{0.3cm}
  {\normalsize
    Street, J., Ihianle, I.K., Olajide, F., \& Lotfi, A. (2025).\\
    \textit{Enhanced Online Grooming Detection Employing Context Determination}\\
    \textit{and Message-Level Analysis.}\\
    \textit{Intelligent Systems with Applications}, 28, 200607.\\
    \href{https://doi.org/10.1016/j.iswa.2025.200607}
         {\texttt{doi:10.1016/j.iswa.2025.200607}}
  \par}

  \vspace{0.8cm}
  \rule{\linewidth}{0.5pt}
  \vspace{0.8cm}

  {\large\textbf{Dataset:}\par}
  \vspace{0.2cm}
  {\normalsize
    Koutli, E. (2023). \textit{PervertedJusticeDataset}.\\
    \href{https://github.com/helenkoutli/PervertedJusticeDataset}
         {\texttt{github.com/helenkoutli/PervertedJusticeDataset}}
  \par}

  \vspace{1.2cm}

  {\large\textbf{Repository:}\par}
  \vspace{0.2cm}
  {\normalsize
    \href{https://github.com/chptt/GroomingDetection}
         {\texttt{github.com/chptt/GroomingDetection}}
  \par}

  \vfill

  {\normalsize BTP-demo Project\par}
  \vspace{0.3cm}
  {\normalsize\today\par}
\end{titlepage}

% ══ Table of contents ═════════════════════════════════════════════════════════
\tableofcontents
\newpage

% ══════════════════════════════════════════════════════════════════════════════
\section{Project Overview}
% ══════════════════════════════════════════════════════════════════════════════

This project implements the online grooming detection system described in
Street et al.\ (2025) using the publicly available PervertedJustice (PJ)
dataset. The system performs two complementary analyses:

\begin{enumerate}[leftmargin=*, label=\textbf{\arabic*.}]
  \item \textbf{Message-Level Analysis (MLA):} Each chat utterance is
        classified as \emph{Adult} (label 1) or \emph{Child} (label 0) based
        on language patterns, using four classifiers: SVM, Naive Bayes, BERT,
        and RoBERTa.

  \item \textbf{Context Determination:} Per-actor predictions are aggregated
        via a binomial significance test to determine whether a full
        conversation involves an Adult--Child (AC) dyad, which is a key signal
        for Online Grooming (OG).
\end{enumerate}

The implementation is faithful to the paper wherever the dataset permits.
All deviations from the paper's methodology are explicitly documented with
justification.

\subsection{Key Clarifications}

\begin{cautionbox}
\textbf{Important:} The following distinctions are enforced throughout the
entire codebase and must not be conflated.
\end{cautionbox}

\vspace{0.5em}
\begin{tabularx}{\linewidth}{>{\itshape}lX}
  \toprule
  \textbf{Common Misconception} & \textbf{Correct Interpretation} \\
  \midrule
  ``Predator utterances are grooming messages'' &
    Predator utterances carry the label \textit{Adult} based on file
    structure, not message content. \\
  \addlinespace
  ``TP\% = grooming detection accuracy'' &
    TP\% measures speaker-role classification accuracy
    (Adult vs.\ Child language style). \\
  \addlinespace
  ``We can detect grooming in new chats'' &
    This system classifies language style; it does not flag
    individual sentences as grooming. \\
  \addlinespace
  ``FP=0 means perfect classification'' &
    FP=0 is a dataset property --- all PJ conversations are OG-positive,
    so no false positive can arise on this dataset alone. \\
  \bottomrule
\end{tabularx}

% ══════════════════════════════════════════════════════════════════════════════
\section{Dataset Acquisition and Inspection}
% ══════════════════════════════════════════════════════════════════════════════

\subsection{Repository}

The dataset was cloned directly from GitHub into the project workspace:

\begin{lstlisting}
git clone https://github.com/helenkoutli/PervertedJusticeDataset.git
\end{lstlisting}

\subsection{Repository Structure}

\begin{lstlisting}
PervertedJusticeDataset/
|-- predators/      # 621 .txt files, one per predator username
|-- officers/       # 638 .txt files, one per officer username
|-- meta.tsv        # 623 rows: name, date, place, comment, URLs
`-- url_name.tsv    # 1244 rows: URL <-> HTML filename mapping
\end{lstlisting}

\subsection{Inspection Findings}

\begin{table}[H]
  \centering
  \caption{PervertedJustice Dataset Statistics}
  \begin{tabularx}{\linewidth}{Xr}
    \toprule
    \textbf{Property} & \textbf{Value} \\
    \midrule
    Predator files                            & 621 \\
    Officer files                             & 638 \\
    Paired conversations (both files exist)   & 609 \\
    Predator-only (no matching officer)       & 12  \\
    Officer-only (multi-officer, \_1/\_2 suffixes) & 29 \\
    Total predator utterances                 & 706,214 \\
    Total officer utterances                  & 662,925 \\
    Grand total utterances                    & 1,369,139 \\
    Utterances per file --- min / max / mean  & 8 / 27,543 / 2,106 \\
    Text encoding                             & ASCII \\
    Empty files                               & 0 \\
    Duplicate files                           & 0 \\
    \bottomrule
  \end{tabularx}
  \label{tab:dataset_stats}
\end{table}

\subsection{File Naming Convention}

\begin{itemize}[leftmargin=*]
  \item \texttt{predators/predator\_\textlangle username\textrangle.txt}
        --- all utterances by the predator.
  \item \texttt{officers/officer\_\textlangle username\textrangle.txt}
        --- all utterances by the officer(s) in that conversation.
  \item Conversations with multiple officers use suffixes:
        \texttt{officer\_\textlangle username\textrangle\_1.txt},
        \texttt{officer\_\textlangle username\textrangle\_2.txt}.
  \item Conversation ID $=$ \textlangle username\textrangle\ extracted
        from the filename.
\end{itemize}

\subsection{Critical Dataset Constraints}

\begin{itemize}[leftmargin=*]
  \item \textbf{Timestamps removed} during cleaning $\Rightarrow$ original
        chronological interleaving is lost; phase-based analysis cannot be
        reproduced.
  \item \textbf{Usernames removed} $\Rightarrow$ actors cannot be
        re-identified.
  \item \textbf{PJ comments / interventions removed} $\Rightarrow$ clean
        utterance-only format.
  \item \textbf{All conversations are OG-positive} $\Rightarrow$ every
        transcript is a confirmed grooming interaction; no negative examples
        exist in this repository.
  \item \textbf{Officers are adult volunteers} posing as children, not real
        minors.
\end{itemize}

% ══════════════════════════════════════════════════════════════════════════════
\section{Paper Methodology Summary}
% ══════════════════════════════════════════════════════════════════════════════

\subsection{Message-Level Analysis (MLA)}

Each utterance is independently classified as Adult (1) or Child (0).

\subsubsection{Models}

\begin{itemize}[leftmargin=*]
  \item \textbf{SVM} --- \texttt{sklearn} Pipeline with
        \texttt{CountVectorizer} + \texttt{LinearSVC} (calibrated).
  \item \textbf{Naive Bayes} --- \texttt{sklearn} \texttt{MultinomialNB}
        with \texttt{CountVectorizer}.
  \item \textbf{BERT} --- \texttt{bert-base-uncased} fine-tuned with
        TensorFlow (3 epochs, learning rate $2 \times 10^{-5}$).
  \item \textbf{RoBERTa} --- \texttt{roberta-base} fine-tuned with
        TensorFlow (3 epochs, learning rate $2 \times 10^{-5}$).
\end{itemize}

\subsubsection{Message Threshold \texorpdfstring{$t$}{t} (Equation 1)}

A message with prediction probability $n$ (probability of Adult class) is
\emph{omitted} when:

\begin{equation}
  (0.5 + t) > n > (0.5 - t)
  \label{eq:threshold}
\end{equation}

Messages falling in this low-confidence band are excluded from metric
calculation. Values of $t$ tested:
$t \in \{0.2,\, 0.3,\, 0.4,\, 0.45\}$.
This applies to BERT and RoBERTa only; SVM and NB produce no omissions.

\subsubsection{MLA Metrics}

\begin{itemize}[leftmargin=*]
  \item \textbf{TP\%} --- correct Adult/Child determinations
        (percentage of \emph{kept} messages).
  \item \textbf{IA\%} --- Incorrect Adult: predicted Adult, true label was
        Child.
  \item \textbf{IC\%} --- Incorrect Child: predicted Child, true label was
        Adult.
  \item \textbf{O\%} --- Omissions: messages excluded by threshold $t$
        (percentage of \emph{all} messages, not just kept ones).
\end{itemize}

Note: TP\% $+$ IA\% $+$ IC\% $\approx$ 100\% over kept messages; O\%
is reported separately over all messages.

\subsection{Context Determination}

For each conversation, both actors (predator and officer) are evaluated
independently using the following algorithm (paper Figure 3):

\begin{enumerate}[leftmargin=*, label=\textbf{Step \arabic*.}]
  \item Collect per-actor $P(\text{Adult})$ predictions from the MLA model.
  \item Apply threshold $t$ (Equation~\ref{eq:threshold}) to omit
        low-confidence messages.
  \item Run a two-sided binomial test:
        \[
          H_0: P(\text{Adult vote}) = 0.5
        \]
        using \texttt{scipy.stats.binomtest}.
  \item Apply the \textbf{Actor Significance Threshold (AST)}:
        \begin{align*}
          p\text{-value} < \text{AST}\ \&\ \bar{n} > 0.5
            &\;\Rightarrow\; \text{Actor} = \textbf{Adult (A)} \\
          p\text{-value} < \text{AST}\ \&\ \bar{n} \leq 0.5
            &\;\Rightarrow\; \text{Actor} = \textbf{Child (C)} \\
          p\text{-value} \geq \text{AST}
            &\;\Rightarrow\; \text{Actor} = \textbf{NS (Not Significant)}
        \end{align*}
        AST values tested: $\{0.05,\, 0.01,\, 0.001\}$.
  \item Apply the \textbf{AC Context Rule} (paper Figure 3):
        \begin{equation}
          \text{is\_AC} =
            \bigl(A \oplus B \text{ is Adult}\bigr)
            \;\wedge\;
            \bigl(A \oplus B \text{ is Child}\bigr)
          \label{eq:ac_rule}
        \end{equation}
        i.e., exactly one actor is Adult \emph{and} exactly one is Child;
        NS in either position $\Rightarrow$ non-AC.
\end{enumerate}

Context metrics: TP, FP, TN, FN, F1-score.

\subsection{Train / Test Split}

80\% training / 20\% testing, split \textbf{by conversation} (not by
individual utterance). The same conversation's predator and officer files
always land in the same split, preventing any data leakage.
Implemented using \texttt{sklearn.model\_selection.GroupShuffleSplit}.

\subsection{Preprocessing}

Per paper Section~4.2: \textbf{no preprocessing} beyond excluding
transcripts shorter than 10 messages (for the PAN12 dataset). Punctuation,
stopwords, and casing are preserved deliberately --- they carry signal for
Adult/Child language distinction.

% ══════════════════════════════════════════════════════════════════════════════
\section{Paper-to-Dataset Reproducibility Mapping}
% ══════════════════════════════════════════════════════════════════════════════

\begin{longtable}{>{\raggedright\arraybackslash}p{4.2cm}
                  >{\centering\arraybackslash}p{2.1cm}
                  >{\raggedright\arraybackslash}p{7.5cm}}
  \caption{Reproducibility of Each Paper Component}
  \label{tab:reproducibility} \\
  \toprule
  \textbf{Component} & \textbf{Status} & \textbf{Notes} \\
  \midrule
  \endfirsthead
  \multicolumn{3}{c}{\tablename~\thetable\ (continued)} \\
  \toprule
  \textbf{Component} & \textbf{Status} & \textbf{Notes} \\
  \midrule
  \endhead
  \bottomrule
  \endfoot

  MLA --- SVM                & \approxbadge  &
    Reproducible; lower TP on small samples due to training-set size. \\
  \addlinespace
  MLA --- Naive Bayes        & \approxbadge  &
    Same as SVM. \\
  \addlinespace
  MLA --- BERT               & \approxbadge  &
    Requires TensorFlow + GPU; fully implemented. \\
  \addlinespace
  MLA --- RoBERTa            & \approxbadge  &
    Requires TensorFlow + GPU; fully implemented. \\
  \addlinespace
  Message Threshold $t$ (Eq.~\ref{eq:threshold}) & \exactbadge &
    Implemented precisely as stated in the paper. \\
  \addlinespace
  Actor Significance (binomial test) & \exactbadge &
    \texttt{scipy.stats.binomtest} with two-sided alternative. \\
  \addlinespace
  AC Context Rule (Eq.~\ref{eq:ac_rule}) & \exactbadge &
    XOR logic implemented exactly as described. \\
  \addlinespace
  MLA metrics (TP, IA, IC, O) & \exactbadge &
    Matches paper Tables~2/3 format exactly. \\
  \addlinespace
  Context metrics (TP, FP, TN, FN, F1) & \partialbadge &
    FP=0, TN=0 on PJ (all-positive dataset); F1 valid for TP/FN. \\
  \addlinespace
  80/20 conversation split   & \exactbadge  &
    \texttt{GroupShuffleSplit} by \texttt{conversation\_id}. \\
  \addlinespace
  Minimal preprocessing      & \exactbadge  &
    Whitespace strip only; no punctuation or stopword removal. \\
  \addlinespace
  Grooming phase analysis    & \notbadge    &
    Timestamps removed; chronological order is lost. \\
  \addlinespace
  Cross-dataset eval (PAN12) & \notbadge    &
    PAN12 not included; requires a separate download. \\
  \addlinespace
  Turn-taking / response time & \notbadge   &
    Speakers separated into different files; interleaving lost. \\
\end{longtable}

\textbf{Legend:}
\exactbadge\ Fully reproducible as described \quad
\approxbadge\ Minor differences expected \quad
\partialbadge\ Dataset-imposed limitation \quad
\notbadge\ Not possible / not reproduced

% ══════════════════════════════════════════════════════════════════════════════
\section{Implementation Architecture}
% ══════════════════════════════════════════════════════════════════════════════

\subsection{Directory Layout}

\begin{lstlisting}
BTP-demo/
|-- PervertedJusticeDataset/        # Cloned dataset (never modified)
|   |-- predators/   (621 files)
|   |-- officers/    (638 files)
|   |-- meta.tsv
|   `-- url_name.tsv
|
`-- grooming_detection/             # All implementation code
    |-- inspect_dataset.py
    |-- data_loader.py
    |-- preprocessing.py
    |-- message_classifier.py
    |-- context_determination.py
    |-- actor_classification.py
    |-- evaluation.py
    |-- run_experiments.py
    |-- requirements.txt
    |-- README.md
    |-- REPORT.md
    |-- ui/
    |   |-- app.py                  # Streamlit entry point
    |   `-- pages/
    |       |-- home.py
    |       |-- dataset_overview.py
    |       |-- run_experiments_ui.py
    |       |-- results_viewer.py
    |       |-- live_classifier.py
    |       `-- paper_mapping.py
    |-- tests/
    |   |-- test_file_pairing.py
    |   |-- test_role_labels.py
    |   |-- test_preprocessing.py
    |   |-- test_data_loader.py
    |   `-- test_splitting.py
    `-- results/
        |-- dataset_inspection_report.txt
        |-- paper_dataset_mapping.md
        |-- mla_svm.json / mla_nb.json
        |-- context_svm.csv / context_nb.csv
        |-- confusion_matrix_svm.png / confusion_matrix_nb.png
        `-- experiment_summary.json
\end{lstlisting}

\subsection{Data Flow}

\begin{center}
\begin{tabular}{ccccc}
  \texttt{data\_loader.py}
  & $\longrightarrow$
  & \texttt{preprocessing.py}
  & $\longrightarrow$
  & \texttt{message\_classifier.py} \\
  \multicolumn{5}{c}{$\downarrow$} \\
  \multicolumn{5}{c}{\texttt{actor\_classification.py}} \\
  \multicolumn{5}{c}{$\downarrow$} \\
  \multicolumn{5}{c}{\texttt{context\_determination.py}
    $\longrightarrow$ \texttt{evaluation.py}} \\
\end{tabular}
\end{center}

% ══════════════════════════════════════════════════════════════════════════════
\section{Module Descriptions}
% ══════════════════════════════════════════════════════════════════════════════

\subsection{\texttt{inspect\_dataset.py}}

Standalone script that analyses the raw dataset and produces
\texttt{results/dataset\_inspection\_report.txt}. Reports file counts,
text encoding, utterance statistics, pairing validation, duplicate
detection, and reproducibility assessment.

\subsection{\texttt{data\_loader.py}}

Core data-access layer. Loads all utterances from both speaker folders
into a \texttt{pandas.DataFrame} with the following columns:

\begin{itemize}[leftmargin=*]
  \item \texttt{utterance} --- message text.
  \item \texttt{speaker\_role} --- \texttt{1} (predator / Adult) or
        \texttt{0} (officer / Child); assigned from \textbf{file structure
        only}, never from message content.
  \item \texttt{conversation\_id} --- username extracted from filename;
        groups predator and officer utterances from the same conversation.
  \item \texttt{file\_name} --- source filename for full traceability.
\end{itemize}

\subsection{\texttt{preprocessing.py}}

\begin{itemize}[leftmargin=*]
  \item \texttt{minimal\_clean()} --- whitespace strip only (paper \S4.2).
  \item \texttt{group\_train\_test\_split()} --- 80/20 split by
        conversation using \texttt{GroupShuffleSplit}; guarantees zero
        leakage between splits.
  \item \texttt{build\_count\_vectorizer()} --- fits \texttt{CountVectorizer}
        on training texts (paper's feature-extraction method).
  \item \texttt{extract\_labels\_and\_texts()} --- splits a DataFrame into
        \texttt{(texts, labels)} arrays.
\end{itemize}

\subsection{\texttt{message\_classifier.py}}

Implements the paper's MLA with four classifiers:

\begin{itemize}[leftmargin=*]
  \item \texttt{SVMClassifier} --- \texttt{LinearSVC} with calibration.
  \item \texttt{NaiveBayesClassifier} --- \texttt{MultinomialNB}.
  \item \texttt{BERTClassifier} --- \texttt{bert-base-uncased}
        (TensorFlow / HuggingFace).
  \item \texttt{RoBERTaClassifier} --- \texttt{roberta-base}
        (TensorFlow / HuggingFace).
\end{itemize}

Also contains:
\begin{itemize}[leftmargin=*]
  \item \texttt{apply\_message\_threshold(proba,~t)} --- implements
        Equation~\ref{eq:threshold}.
  \item \texttt{compute\_mla\_metrics(y\_true, y\_pred, \ldots)} ---
        computes TP\%, IA\%, IC\%, O\%.
\end{itemize}

\subsection{\texttt{context\_determination.py}}

Implements the paper's Figure~3 algorithm:

\begin{itemize}[leftmargin=*]
  \item \texttt{determine\_actor\_significance()} --- runs
        \texttt{scipy.stats.binomtest} and applies AST; returns
        \textbf{A}, \textbf{C}, or \textbf{NS}.
  \item \texttt{determine\_transcript\_context()} --- applies the XOR
        logic (Equation~\ref{eq:ac_rule}) to produce AC or non-AC.
  \item \texttt{run\_context\_determination()} --- processes a full test
        DataFrame conversation by conversation.
  \item \texttt{compute\_context\_metrics()} --- computes TP, FP, TN,
        FN, F1.
\end{itemize}

\subsection{\texttt{actor\_classification.py}}

Pipeline bridge between message classifiers and context determination:

\begin{itemize}[leftmargin=*]
  \item \texttt{attach\_predictions()} --- runs a trained classifier over
        all test utterances and attaches $P(\text{Adult})$ probabilities.
  \item \texttt{sweep\_ast\_t\_values()} --- sweeps all combinations of
        AST $\times$ $t$ and returns a results DataFrame.
  \item \texttt{run\_full\_actor\_analysis()} --- runs the full sweep for
        all classifiers.
\end{itemize}

\subsection{\texttt{evaluation.py}}

\begin{itemize}[leftmargin=*]
  \item \texttt{compute\_standard\_metrics()} --- accuracy, precision,
        recall, F1 via \texttt{sklearn}.
  \item \texttt{plot\_confusion\_matrix()} --- saves a labeled confusion
        matrix PNG (labeled as speaker-role classification, not grooming).
  \item \texttt{format\_mla\_table()} / \texttt{format\_context\_table()}
        --- human-readable results tables.
  \item \texttt{print\_paper\_comparison()} --- side-by-side display with
        published results (reference only; never injected into outputs).
\end{itemize}

\subsection{\texttt{run\_experiments.py}}

Full pipeline orchestration. CLI flags:
\begin{itemize}[leftmargin=*]
  \item \texttt{--skip-transformers} --- run SVM and NB only.
  \item \texttt{--fast} --- sample $N$ conversations for quick
        verification.
  \item \texttt{--debug} --- verbose logging.
\end{itemize}

\subsection{\texttt{ui/} --- Streamlit Web Application}

Six-page interactive web UI (Section~\ref{sec:ui}).

% ══════════════════════════════════════════════════════════════════════════════
\section{Experimental Results}
% ══════════════════════════════════════════════════════════════════════════════

\subsection{Fast-Mode Baseline (50-Conversation Sample)}

The following results were obtained using a 50-conversation sample for
quick verification. Full dataset results will differ; see
Section~\ref{sec:paper_results} for the paper's published targets.

\subsubsection{Message-Level Analysis}

\begin{table}[H]
  \centering
  \caption{MLA Results --- Fast Mode (50 conversations, 10-conversation test set)}
  \begin{tabular}{lrrrr}
    \toprule
    \textbf{Model} & \textbf{TP\%} & \textbf{IA\%} & \textbf{IC\%} &
    \textbf{O\%} \\
    \midrule
    SVM & 61.8 & 28.7 & 9.5 & 0.0 \\
    NB  & 57.0 & 35.2 & 7.8 & 0.0 \\
    \bottomrule
  \end{tabular}
  \label{tab:mla_fast}
\end{table}

\begin{infobox}
These are \textbf{speaker-role classification} metrics (Adult vs.\ Child
language style). They are \textbf{not} grooming detection metrics. Lower
TP\% compared to the paper is expected: fast mode trains on approximately
85\,000 utterances, whereas the paper uses over 219\,000.
\end{infobox}

\subsubsection{Context Determination (10 Test Conversations)}

\begin{table}[H]
  \centering
  \caption{Context Determination --- Fast Mode}
  \begin{tabular}{llrrrrr}
    \toprule
    \textbf{Model} & \textbf{AST} & \textbf{TP} & \textbf{FN} &
    \textbf{FP} & \textbf{TN} & \textbf{F1} \\
    \midrule
    SVM & 0.050 & 4 & 3 & 0 & 0 & 0.727 \\
    SVM & 0.010 & 4 & 3 & 0 & 0 & 0.727 \\
    SVM & 0.001 & 4 & 3 & 0 & 0 & 0.727 \\
    \addlinespace
    NB  & 0.050 & 5 & 2 & 0 & 0 & 0.833 \\
    NB  & 0.010 & 5 & 2 & 0 & 0 & 0.833 \\
    NB  & 0.001 & 5 & 2 & 0 & 0 & 0.833 \\
    \bottomrule
  \end{tabular}
  \label{tab:ctx_fast}
\end{table}

FP$=0$ and TN$=0$ are correct and expected because all PJ conversations
are OG-positive; there are no true negative transcripts in this dataset.

\subsection{Paper Published Results (Reference Only)}
\label{sec:paper_results}

\begin{cautionbox}
These values are from Street et al.\ (2025) and are shown for comparison
only. They are never embedded in experiment outputs or used to validate
our runs.
\end{cautionbox}

\subsubsection{MLA --- PJ Dataset, Inter-Set (Paper Table 2)}

\begin{table}[H]
  \centering
  \caption{Paper MLA Results --- PJ Dataset, Inter-Set Training}
  \begin{tabular}{llrrrr}
    \toprule
    \textbf{Model} & \textbf{$t$} & \textbf{TP\%} & \textbf{IA\%} &
    \textbf{IC\%} & \textbf{O\%} \\
    \midrule
    SVM       & ---  & 76.0 & 13.5 & 10.5 &  0.0 \\
    NB        & ---  & 74.7 & 12.0 & 13.3 &  0.0 \\
    BERT      & 0.20 & 82.2 &  9.3 &  8.5 &  9.2 \\
    BERT      & 0.45 & 87.3 &  7.1 &  5.6 & 30.0 \\
    RoBERTa   & 0.20 & 83.3 &  7.7 &  9.0 &  8.8 \\
    RoBERTa   & 0.45 & 88.8 &  5.3 &  5.9 & 30.4 \\
    \bottomrule
  \end{tabular}
  \label{tab:mla_paper}
\end{table}

\subsubsection{Context Determination --- PJ, AST$=0.05$ (Paper Table 4)}

\begin{table}[H]
  \centering
  \caption{Paper Context Determination Results --- PJ Dataset, AST$=0.05$}
  \begin{tabular}{lrrr}
    \toprule
    \textbf{Model} & \textbf{TP} & \textbf{FN} & \textbf{F1} \\
    \midrule
    SVM                   & 118 & 35 & 0.871 \\
    NB                    & 116 & 37 & 0.862 \\
    BERT ($t=0.2$)        & 143 & 10 & 0.966 \\
    RoBERTa ($t=0.45$)    & 145 &  8 & 0.973 \\
    \bottomrule
  \end{tabular}
  \label{tab:ctx_paper}
\end{table}

\subsection{Reasons Our Results Differ from the Paper}

\begin{enumerate}[leftmargin=*, label=\textbf{\arabic*.}]
  \item \textbf{Dataset version:} The paper parsed HTML directly from
        perverted-justice.com; this repository is a pre-cleaned subset
        of 622 conversations.
  \item \textbf{Training set size:} Fast mode uses 50 conversations;
        the paper uses the full dataset.
  \item \textbf{Framework:} The paper uses TensorFlow throughout; SVM
        and NB here use \texttt{sklearn} (consistent with the paper's
        own \texttt{sklearn} reference).
  \item \textbf{Randomness:} Split randomness affects results
        significantly on small samples.
\end{enumerate}

% ══════════════════════════════════════════════════════════════════════════════
\section{Web User Interface}
\label{sec:ui}
% ══════════════════════════════════════════════════════════════════════════════

A six-page Streamlit web application provides interactive access to all
functionality. Start the application with:

\begin{lstlisting}
cd grooming_detection
python -m streamlit run ui/app.py
# Opens at http://localhost:8501
\end{lstlisting}

\begin{table}[H]
  \centering
  \caption{Web UI Pages and Their Functionality}
  \begin{tabularx}{\linewidth}{>{\bfseries}lX}
    \toprule
    Page & Description \\
    \midrule
    Home &
      Status dashboard showing dataset availability, file counts, test
      pass rate, and constraint warnings. \\
    \addlinespace
    Dataset Overview &
      Interactive histograms of utterances per file, role balance pie
      chart, pairing analysis, and a searchable per-file table. \\
    \addlinespace
    Run Experiments &
      Interactive controls: model selection, test-split slider, fast mode,
      AST sweep. Streams live training progress and displays results
      inline. \\
    \addlinespace
    Results \& Metrics &
      MLA grouped bar chart, context determination F1 heatmap (AST
      $\times$ model), confusion matrix images, and side-by-side
      comparison with paper reference results. \\
    \addlinespace
    Live Classifier &
      Real-time message classification with a probability gauge; batch
      classification with a colour-coded table; demo conversation with
      chat-bubble UI. \\
    \addlinespace
    Paper Mapping &
      Full reproducibility table, three-task explanation (MLA vs.\
      actor significance vs.\ context determination), and the raw
      dataset inspection report. \\
    \bottomrule
  \end{tabularx}
  \label{tab:ui_pages}
\end{table}

% ══════════════════════════════════════════════════════════════════════════════
\section{Test Suite}
% ══════════════════════════════════════════════════════════════════════════════

\textbf{88 tests, all passing.}

\begin{lstlisting}
python -m pytest grooming_detection/tests/ -v
# 88 passed in 13.60s
\end{lstlisting}

\begin{table}[H]
  \centering
  \caption{Test Files and Coverage}
  \begin{tabularx}{\linewidth}{lrX}
    \toprule
    \textbf{File} & \textbf{Tests} & \textbf{What is Verified} \\
    \midrule
    \texttt{test\_file\_pairing.py}   & 14 &
      Filename parsing, predator/officer file counts, multi-officer
      suffix handling, no duplicates, no empty files. \\
    \addlinespace
    \texttt{test\_role\_labels.py}    &  9 &
      Labels assigned from file structure not content; both roles
      present; no \texttt{is\_grooming} column. \\
    \addlinespace
    \texttt{test\_preprocessing.py}  & 18 &
      Minimal cleaning preserves punctuation and casing; group split
      has zero conversation leakage; reproducibility with fixed seed. \\
    \addlinespace
    \texttt{test\_data\_loader.py}   & 11 &
      Temp-directory loading; multi-officer files; min-utterance filter;
      null check; conversation count. \\
    \addlinespace
    \texttt{test\_splitting.py}      & 20 &
      Group integrity; message threshold boundary (Eq.~\ref{eq:threshold});
      MLA metric correctness; binomial test direction; AC/non-AC XOR logic. \\
    \midrule
    \textbf{Total} & \textbf{72} & \emph{Unit tests (no dataset required)} \\
    & \textbf{16} & \emph{Integration tests (dataset required)} \\
    \bottomrule
  \end{tabularx}
  \label{tab:tests}
\end{table}

\subsection{Key Safety Tests}

\begin{itemize}[leftmargin=*]
  \item \textbf{No conversation leakage} --- \texttt{test\_zero\_conversation\_overlap}
        verifies train/test conversation IDs are disjoint.
  \item \textbf{Label source} ---
        \texttt{test\_label\_source\_is\_file\_structure\_not\_content}
        verifies labels are derived from folder membership, not from message
        text.
  \item \textbf{No grooming column} ---
        \texttt{test\_label\_docstring\_warning\_not\_violated} asserts
        no \texttt{is\_grooming} column is created.
  \item \textbf{Threshold boundaries} ---
        \texttt{test\_boundary\_values\_not\_omitted} verifies the strict
        inequality in Equation~\ref{eq:threshold}.
  \item \textbf{XOR logic} ---
        \texttt{test\_both\_adult\_not\_ac} and
        \texttt{test\_ns\_in\_either\_position\_not\_ac} verify that the
        context rule correctly rejects invalid configurations.
\end{itemize}

% ══════════════════════════════════════════════════════════════════════════════
\section{Known Limitations}
% ══════════════════════════════════════════════════════════════════════════════

\subsection{Dataset Limitations}

\begin{table}[H]
  \centering
  \caption{Dataset-Level Limitations}
  \begin{tabularx}{\linewidth}{>{\raggedright}p{3.8cm}
                               >{\raggedright}p{4cm}
                               X}
    \toprule
    \textbf{Limitation} & \textbf{Impact} & \textbf{Path to Resolve} \\
    \midrule
    Timestamps removed &
      Phase-based analysis and early detection cannot be reproduced. &
      Document and skip those paper components. \\
    \addlinespace
    All-positive dataset &
      FP$=0$, TN$=0$ --- false positive rate cannot be measured. &
      Add PAN12 dataset for cross-dataset evaluation. \\
    \addlinespace
    Officers are adult volunteers &
      Child language model may not generalise to real children. &
      Acknowledge in all reporting. \\
    \addlinespace
    622 conversations only &
      Smaller than the full PJ corpus used in the paper. &
      Use the full dataset for final runs. \\
    \bottomrule
  \end{tabularx}
  \label{tab:dataset_limits}
\end{table}

\subsection{Implementation Limitations}

\begin{table}[H]
  \centering
  \caption{Implementation-Level Limitations}
  \begin{tabularx}{\linewidth}{>{\raggedright}p{4.5cm}
                               >{\raggedright}p{3.5cm}
                               X}
    \toprule
    \textbf{Limitation} & \textbf{Impact} & \textbf{Path to Resolve} \\
    \midrule
    BERT/RoBERTa require TensorFlow + GPU &
      Transformer experiments not runnable on CPU-only machines. &
      Install \texttt{tensorflow==2.17.0} and run on GPU. \\
    \addlinespace
    PAN12 not included &
      Cross-dataset F1 scores (paper Tables~5--7) not reproduced. &
      Download PAN12 separately and add \texttt{pan12\_loader.py}. \\
    \addlinespace
    Fast-mode results differ from paper &
      Expected --- not all 622 conversations are used. &
      Run without \texttt{--fast} for full results. \\
    \bottomrule
  \end{tabularx}
  \label{tab:impl_limits}
\end{table}

% ══════════════════════════════════════════════════════════════════════════════
\section{Reproduction Instructions}
% ══════════════════════════════════════════════════════════════════════════════

\subsection{Prerequisites}

\begin{lstlisting}
# Python 3.10+
pip install -r grooming_detection/requirements.txt

# For BERT/RoBERTa (optional, GPU recommended):
pip install tensorflow==2.17.0
\end{lstlisting}

\subsection{Step-by-Step}

\begin{lstlisting}
# 1. Clone dataset
git clone https://github.com/helenkoutli/PervertedJusticeDataset.git

# 2. Inspect dataset
python grooming_detection/inspect_dataset.py

# 3. Run test suite
python -m pytest grooming_detection/tests/ -v

# 4. Quick verification (50 conversations, ~30 seconds)
python grooming_detection/run_experiments.py --skip-transformers --fast

# 5. Full SVM + NB on complete dataset
python grooming_detection/run_experiments.py --skip-transformers

# 6. Full pipeline with BERT + RoBERTa (GPU required)
python grooming_detection/run_experiments.py

# 7. Launch web UI
python -m streamlit run grooming_detection/ui/app.py
\end{lstlisting}

\subsection{Expected Outputs}

After running Step~5 (full SVM + NB):

\begin{lstlisting}
results/
|-- dataset_inspection_report.txt   # dataset statistics
|-- paper_dataset_mapping.md        # reproducibility table
|-- mla_svm.json                    # SVM MLA metrics
|-- mla_nb.json                     # NB MLA metrics
|-- context_svm.csv                 # SVM context determination sweep
|-- context_nb.csv                  # NB context determination sweep
|-- confusion_matrix_svm.png        # SVM confusion matrix
|-- confusion_matrix_nb.png         # NB confusion matrix
`-- experiment_summary.json         # combined summary
\end{lstlisting}

% ══════════════════════════════════════════════════════════════════════════════
\section*{References}
\addcontentsline{toc}{section}{References}
% ══════════════════════════════════════════════════════════════════════════════

\begin{enumerate}[leftmargin=*, label={[\arabic*]}]

  \item Street, J., Ihianle, I.K., Olajide, F., \& Lotfi, A. (2025).
        Enhanced online grooming detection employing context determination
        and message-level analysis.
        \textit{Intelligent Systems with Applications}, 28, 200607.
        \href{https://doi.org/10.1016/j.iswa.2025.200607}
             {\texttt{doi:10.1016/j.iswa.2025.200607}}

  \item Koutli, E. (2023). PervertedJusticeDataset [Dataset].
        GitHub.
        \href{https://github.com/helenkoutli/PervertedJusticeDataset}
             {\texttt{github.com/helenkoutli/PervertedJusticeDataset}}

  \item Pedregosa, F. et al.\ (2011). Scikit-learn: Machine learning in
        Python. \textit{Journal of Machine Learning Research}, 12,
        2825--2830.

  \item Devlin, J., Chang, M.W., Lee, K., \& Toutanova, K. (2019).
        BERT: Pre-training of deep bidirectional transformers for language
        understanding. \textit{NAACL-HLT 2019}, 4171--4186.

  \item Liu, Y. et al.\ (2019). RoBERTa: A robustly optimised BERT
        pretraining approach. \textit{arXiv:1907.11692}.

  \item Perverted Justice Foundation.
        \href{http://www.perverted-justice.com}
             {\texttt{www.perverted-justice.com}}

\end{enumerate}

% ══════════════════════════════════════════════════════════════════════════════
\vfill
\begin{center}
  \small\textit{Report generated for BTP-demo project.}\\
  \small\textit{Dataset used for academic research only.}
\end{center}

\end{document}
